\section{Explication vulgarisée originale du LLM}
\label{app:E}

\textbf{Texte du LLM reproduit fidèlement.} La section B ci-dessous est extraite de \sourcepath{llm/responses/phase8_official_llm_response.md}, sans correction de son texte, de ses valeurs ni de ses références. Elle est incluse selon l’exigence de la section 5 du cahier des charges transmise par l’équipe. L’audit critique demeure au chapitre 10. La source officielle reste inchangée, SHA-256 : \sourcehash{eecb0de3e0dd8e152e93db70ff71f9dad87024b7b2f3e1bee1ede05225f6f3b4} ; l’empreinte de l’extrait est consignée dans \sourcepath{docs/phase9b_results_traceability.csv}.

\begin{quote}\small

B. VULGARISATION

La stationnarité revient à demander si une série conserve à peu près les mêmes propriétés statistiques au cours du temps. Un prix boursier peut monter pendant plusieurs années et changer fortement de niveau ; il est alors difficile de le traiter comme une série qui oscille autour d’un comportement stable. Dans ce projet, les tests ADF et KPSS conduisent ensemble à traiter les prix et log-prix comme non stationnaires, tandis que les log-rendements sont compatibles avec la stationnarité. ADF et KPSS ne démontrent toutefois pas une vérité absolue : chacun teste une hypothèse nulle différente, et « ne pas rejeter » n’est pas « prouver ».

Le benchmark Naive ou random walk est extrêmement simple : pour prévoir le prix futur, on conserve le dernier prix observé. Ce modèle ne cherche presque aucune structure supplémentaire. Sa simplicité constitue justement un test exigeant : si un modèle complexe ne le bat pas de manière robuste, sa complexité n’a pas apporté de valeur prédictive démontrée.

ARIMA cherche à exploiter la dépendance temporelle d’une série après les transformations nécessaires. Dans ce projet, il était raisonnable d’utiliser une différenciation d’ordre 1 du log-prix. Pourtant, lorsque plusieurs ARIMA ont été comparés en validation interne, aucun challenger n’a satisfait la règle complète permettant de remplacer Naive. Cela illustre une distinction essentielle : bien expliquer le passé et bien prévoir le futur ne sont pas la même chose.

VAR élargit l’analyse à plusieurs séries. Au lieu de regarder seulement le passé de NVDA, on regarde aussi les historiques d’AMD et SOXX. Les tests de Granger indiquent ici que le passé de SOXX ajoute de l’information prédictive pour NVDA et AMD. Le mot « causalité » dans « causalité de Granger » est trompeur pour un non-spécialiste : il s’agit d’antériorité prédictive conditionnelle, pas d’une preuve économique selon laquelle une variation de SOXX produirait mécaniquement une variation de NVDA.

La cointégration cherche à savoir si plusieurs séries non stationnaires partagent malgré tout une relation de long terme suffisamment stable. Ici, Johansen retient un rang nul dans les spécifications testées. La conclusion correcte est donc : aucune relation de cointégration n’est statistiquement soutenue dans ces spécifications. Il serait excessif de dire qu’aucune relation économique de long terme ne peut exister.

GARCH répond encore à une autre question. ARIMA et VAR s’intéressent surtout à la dynamique de moyenne ; GARCH s’intéresse à la variance conditionnelle, donc à l’intensité des fluctuations. Les marchés alternent souvent périodes calmes et périodes agitées : c’est le « volatility clustering ». Ici, les ARCH-LM avant GARCH détectent cette structure, puis les diagnostics après GARCH n’en détectent plus aux horizons testés. La volatilité peut donc être modélisée séparément de la direction du prix.

La prévision GARCH de sigma à l’horizon 10 vaut environ 3.17 \% par jour pour le modèle Normal verrouillé à partir du 14/11/2025. Cela ne veut pas dire que NVDA doit monter ou baisser de 3.17 \%. Sigma mesure l’ampleur conditionnelle des fluctuations, pas leur signe. Source : outputs/tables/phase5\_volatility\_forecast\_train\_end.csv.

Les réseaux LSTM et GRU peuvent apprendre des relations non linéaires à partir de fenêtres historiques. Pourtant, ils perdent ici contre Naive. Plusieurs mécanismes sont plausibles --- bruit d’estimation, quantité de données limitée, extrapolation d’un niveau de log-prix non stationnaire, changement de régime --- mais le projet ne démontre pas lequel explique effectivement l’écart. Ce sont donc des limites plausibles, pas des causes établies.

RMSE et MAE mesurent toutes deux l’erreur de prévision en USD ajustés. RMSE pénalise davantage les grosses erreurs parce qu’elle élève les erreurs au carré avant agrégation. MAE mesure l’écart absolu moyen et est moins dominée par les grandes erreurs. Sur le holdout final à h=5, Naive a RMSE 9.68557 et MAE 7.68782 ; à h=10, RMSE 11.4858 et MAE 9.07453. Source : outputs/tables/phase7\_final\_metrics.csv.

Il faut enfin distinguer validation interne et holdout final. Les Phases 3B et 6 servent à choisir ou verrouiller les spécifications à partir de TRAIN. La Phase 7 est l’évaluation finale sur une période mise à part pour la comparaison hors échantillon. Elle utilise 211 origines communes et exactement les mêmes dates cibles pour tous les modèles. Les résultats de cette phase ne doivent plus servir à retuner les modèles.
\end{quote}
\sourcecaption{Une copie UTF-8 exacte de cet extrait est conservée dans \sourcepath{appendices/llm_section_b_original.txt} ; son SHA256 est dans \sourcepath{sources_manifest.json}. Seule la composition typographique change.}
